% Tabla: Validación de K Vecinos para Cálculos Internos
% Experimento 03 - K-Fold CV (5×3 = 15 folds)
\begin{table}[htbp]
\centering
\caption{Validación de K vecinos para ponderación y filtrado}
\label{tab:k_neighbors_validation}
\begin{tabular}{cccccl}
\toprule
\textbf{K} & \textbf{Macro F1} & \textbf{$\Delta$ F1 (pp)} & \textbf{p-valor} & \textbf{Tasa Acept.} & \textbf{Contexto} \\
\midrule
5   & 0.2063 & +0.18 & 0.042 & 99\% & Insuficiente \\
8   & 0.2064 & +0.19 & 0.038 & 96\% & Limitado \\
10  & 0.2056 & +0.11 & 0.089 & 95\% & Limitado \\
12  & 0.2055 & +0.10 & 0.095 & 98\% & Óptimo \\
15  & 0.2059 & +0.14 & 0.061 & 96\% & Óptimo \\
18  & 0.2065 & +0.20 & 0.028 & 96\% & Óptimo \\
20  & 0.2060 & +0.15 & 0.054 & 97\% & Óptimo \\
25  & 0.2062 & +0.17 & 0.045 & 96\% & Redundante \\
50  & 0.2046 & +0.01 & 0.621 & 97\% & Redundante \\
100 & 0.2063 & +0.18 & 0.043 & 96\% & Ruidoso \\
\rowcolor{green!20}
\textbf{200} & \textbf{0.2078} & \textbf{+0.33} & \textbf{0.0085} & \textbf{97\%} & \textbf{Amplio} \\
\bottomrule
\end{tabular}
\vspace{0.5em}\footnotesize

\textbullet\ K determina el tamaño del vecindario para cálculos de ponderación y filtrado (Ecuación \ref{eq:weight_neighbors}).

\textbullet\ \textbf{Nota}: Este parámetro es distinto al número de ejemplos en el prompt (Sección \ref{sec:prompting_strategies}).

\textbullet\ Línea base Macro F1: 0.2045. Mejor configuración: K=200 (+0.33 pp, p=0.0085).

\end{table}
